\begin{theorem}
Let $E$ be a finite set.  Every super-sequence $f:F\to E$ admits a constant
sub-super-sequence.
\end{theorem}
\begin{proof}
We induct on the finite cardinality of $E$.  For a two-part split, apply
\noderef{NashWilliams} to the subset
$S=\{s\in F\mid f(s)=e\}$ for a chosen color $e$.  If the resulting
sub-front is contained in $S$, the restriction of $f$ to it is constantly
$e$.  If it is disjoint from $S$, its values lie in the remaining finite
color set, so the induction hypothesis applies to this sub-super-sequence.

Formally, the sub-front supplied by \noderef{NashWilliams} carries its own
front structure and is contained in $F$.  Restrict the value map of
\noderef{SuperSequence} along this inclusion.  In the second case the
restricted map has codomain $E\setminus\{e\}$, whose cardinality is one
smaller; the induction hypothesis gives a further front on which the map is
constant.  Composing the two inclusions gives a sub-front of the original
front and preserves constancy.  The empty finite-color case cannot occur for
a super-sequence, since its domain front is nonempty by density.  Therefore
the induction terminates with a constant sub-super-sequence.
\end{proof}
